\section*{Chapter \ref{ch:b1:field-lab-methods-1} --- Field and Lab Methods I}

\begin{solution}{exo:b1:field-lab-methods-1:1}
At 1:10\,000, \qty{4.9}{m} becomes $4.9/10\,000 = \qty{0.49}{mm}$; at
1:25\,000, $4.9/25\,000 = \qty{0.20}{mm}$, about the width of a pencil line.
\end{solution}

\begin{solution}{exo:b1:field-lab-methods-1:2}
At Reykjavik $D = \ang{-11.28}$: $52 - 11.28 = \ang{040.7}$. At Wellington
$D = \ang{+23.23}$: $52 + 23.23 = \ang{075.2}$. The same compass reading means
directions \ang{34.5} apart.
\end{solution}

\begin{solution}{exo:b1:field-lab-methods-1:3}
$20\to015$. With $t = \ang{15}$ and $p = \ang{20}$, the downward unit vector is
\[
  (\cos\ang{20}\sin\ang{15},\ \cos\ang{20}\cos\ang{15},\ -\sin\ang{20}) =
  (0.2432,\ 0.9077,\ -0.3420).
\]
\end{solution}

\begin{solution}{exo:b1:field-lab-methods-1:4}
The arithmetic mean is $(350 + 358 + 6)/3 = \ang{238}$, nearly opposite to
all three readings. The horizontal unit vectors $(\sin a, \cos a)$ add up to
$(-0.1736 - 0.0349 + 0.1045,\ 0.9848 + 0.9994 + 0.9945) = (-0.1040, 2.9787)$,
whose \omterm{def:b1:spherical-astronomy:horizontal}{azimuth} is $\operatorname{atan2}(-0.1040, 2.9787) = \ang{-2.0}$, that
is \ang{358.0}: the right mean, in the middle of the readings.
\end{solution}

\begin{solution}{exo:b1:field-lab-methods-1:5}
Mean $175/5 = \ang{35.0}$. Deviations $-1, 2, 0, -2, 1$; sum of squares 10;
$s = \sqrt{10/4} = \ang{1.58}$. Standard error $1.58/\sqrt5 = \ang{0.71}$.
\end{solution}

\begin{solution}{exo:b1:field-lab-methods-1:6}
The normals are $(\sin\ang{30}, 0, \cos\ang{30}) = (0.5, 0, 0.866)$ and
$(-0.5, 0, 0.866)$; $\vec S = (0, 0, 1.732)$ is vertical: the mean is a
horizontal plane, and $\bar R = 1.732/2 = 0.866$, each reading \ang{30} from
the mean. The two readings do not sample one surface with errors: they belong
to two limbs of a \omterm{def:b1:geological-maps:fold}{fold}, two different populations, and their average
describes neither. Each limb is averaged on its own; together they define the
\omterm{def:b1:geological-maps:fold}{fold axis}, the line common to both planes, here horizontal and north--south.
\end{solution}

\begin{solution}{exo:b1:field-lab-methods-1:7}
$\tan p = \tan\ang{12}\cos(\ang{80} - \ang{120}) = 0.2126\times0.7660 =
0.1628$: $p = \ang{9.3}$, the scratches plunging toward \ang{120}. For the
\omterm{def:b1:field-lab-methods-1:trend-plunge}{trend} \ang{170}, $\cos(\ang{-90}) = 0$: the line is horizontal, along the
\omterm{def:b1:geological-maps:attitude}{strike}.
\end{solution}

\begin{solution}{exo:b1:field-lab-methods-1:8}
$0.47/0.20 = 2.35$ years after 1 January 2026, that is in spring 2028. On
1 January 2030: $-0.47 + 4\times0.20 = \ang{+0.33}$, east of geographic north.
\end{solution}

\begin{solution}{exo:b1:field-lab-methods-1:9}
The three normals lie in the vertical plane of \omterm{def:b1:spherical-astronomy:horizontal}{azimuth} \ang{080}, at
\ang{10}, \ang{12} and \ang{14} from the vertical; the outer two are
symmetric about the middle one, so their sum lies along it: the mean is
$080/12$, with $\theta = \ang{2}, 0, \ang{2}$. Then $\bar R = (1 +
2\cos\ang{2})/3 = (1 + 1.998782)/3 = 0.999594$, and $\sqrt{2(1 - \bar R)} =
\sqrt{0.000812} = \qty{0.02850}{rad} = \ang{1.633}$; exactly,
$\theta_{\text{rms}} = \sqrt{(4 + 0 + 4)/3} = \ang{1.633}$.
\end{solution}

\begin{solution}{exo:b1:field-lab-methods-1:10}
The section is a circle of radius $\sqrt{r^2 - h^2}$. Its mean is
$\frac1r\int_0^r\sqrt{r^2 - h^2}\,\mathrm dh = \frac1r\cdot\frac{\pi r^2}{4}
= \frac{\pi r}{4}$, since the integral is the area of a quarter disc. Grains
\qty{0.50}{mm} across appear on average $0.50\times\pi/4 = \qty{0.39}{mm}$
across: a \omterm{def:b1:field-lab-methods-1:thin-section}{thin section} underestimates grain sizes.
\end{solution}

\begin{solution}{exo:b1:field-lab-methods-1:11}
The rotation by $D$ about the vertical,
$(x, y, z) \mapsto (x\cos D + y\sin D,\ -x\sin D + y\cos D,\ z)$, sends
$(\sin\delta\sin\alpha, \sin\delta\cos\alpha, \cos\delta)$ to
$(\sin\delta\sin(\alpha + D), \sin\delta\cos(\alpha + D), \cos\delta)$: it
adds $D$ to every \omterm{def:b1:geological-maps:attitude}{dip direction}. Being linear, it sends the resultant
$\vec S$ to the resultant of the rotated normals; it keeps lengths and the
vertical component, so the mean \omterm{def:b1:geological-maps:attitude}{dip} $\arccos(S_U/\|\vec S\|)$ is unchanged,
and it turns the horizontal part of $\vec S$ by $D$, so the mean \omterm{def:b1:geological-maps:attitude}{dip direction}
increases by $D$. Correcting each reading for the declination, or only their
mean, gives the same result.
\end{solution}

\begin{solution}{exo:b1:field-lab-methods-1:12}
The scatter along the line of steepest \omterm{def:b1:geological-maps:attitude}{dip} is a scatter of the \omterm{def:b1:geological-maps:attitude}{dip}: the
standard error is $2/\sqrt N \le 0.5$ for $N \ge 16$. Across it, a turn of the
normal by \ang{2} changes the \omterm{def:b1:geological-maps:attitude}{dip direction} by $2/\sin\ang{12} = \ang{9.62}$;
the standard error of the mean \omterm{def:b1:geological-maps:attitude}{dip direction} is $9.62/\sqrt N \le 1$ for
$N \ge 92.5$, that is 93 readings. A gentle \omterm{def:b1:geological-maps:attitude}{dip} is quickly known; its
direction is not.
\end{solution}

\begin{solution}{pb:b1:field-lab-methods-1:1}
\textbf{1.} \omterm{def:b1:geological-maps:attitude}{Dip directions}: $967/12 = \ang{80.6}$; \omterm{def:b1:geological-maps:attitude}{dips}: $145/12 =
\ang{12.1}$.

\textbf{2.} \ang{355} and \ang{005} average to \ang{180}, the opposite
direction. Here all the \omterm{def:b1:geological-maps:attitude}{dip directions} lie between \ang{076} and \ang{085},
far from the jump at \ang{000}, so the arithmetic mean works.

\textbf{3.} $\|\vec n\|^2 = \sin^2\delta(\sin^2\alpha + \cos^2\alpha) +
\cos^2\delta = 1$. The downward \omterm{def:b1:geological-maps:attitude}{dip} line has the direction
$(\cos\delta\sin\alpha, \cos\delta\cos\alpha, -\sin\delta)$, and its dot
product with $\vec n$ is $\sin\delta\cos\delta - \cos\delta\sin\delta = 0$; the
\omterm{def:b1:geological-maps:attitude}{strike} line is horizontal and perpendicular to the \omterm{def:b1:geological-maps:attitude}{dip direction},
$(\cos\alpha, -\sin\alpha, 0)$, and its dot product with $\vec n$ is
$\sin\delta\sin\alpha\cos\alpha - \sin\delta\cos\alpha\sin\alpha = 0$.

\textbf{4.} Reading 11 is the same as reading 2: $(0.2222, 0.0352, 0.9744)$.
Reading 12:
\[
  (\sin\ang{11}\sin\ang{82},\ \sin\ang{11}\cos\ang{82},\ \cos\ang{11}) =
  (0.1890,\ 0.0266,\ 0.9816).
\]

\textbf{5.} $\vec S = (2.4747,\ 0.4082,\ 11.7294)$.

\textbf{6.} $\|\vec S\| = \sqrt{6.1241 + 0.1666 + 137.5788} = 11.99457$, and
$\bar R = 11.99457/12 = 0.999547$.

\textbf{7.} $\bar\alpha = \operatorname{atan2}(2.4747, 0.4082) = \ang{80.6}$
and $\bar\delta = \arccos(11.7294/11.99457) = \arccos 0.977893 = \ang{12.1}$:
the vector mean is $080.6/12.1$, the same as the arithmetic means here.

\textbf{8.} \omterm{def:b1:geological-maps:attitude}{Dip directions}: $s = \ang{2.94}$; \omterm{def:b1:geological-maps:attitude}{dips}: $s = \ang{1.68}$.

\textbf{9.} $\theta_{\text{rms}} \approx \sqrt{2\times0.000453} =
\qty{0.0301}{rad} = \ang{1.72}$.

\textbf{10.} $\sqrt{1.68^2 + (\sin\ang{12.1}\times2.94)^2} = \sqrt{2.82 +
0.38} = \ang{1.79}$, close to \ang{1.72}: the readings behave as the error
model says, within the effects of rounding to the degree and of a small
sample.

\textbf{11.} The \omterm{def:b1:geological-maps:attitude}{dip}: it contributes 2.82 of the 3.20 (in square degrees).
The \omterm{def:b1:geological-maps:attitude}{dip direction} counts only through $\sin\delta = 0.21$, because on a gentle
bed a change of \omterm{def:b1:geological-maps:attitude}{dip direction} turns the plane very little.

\textbf{12.} Mean \omterm{def:b1:geological-maps:attitude}{dip}: $1.68/\sqrt{12} = \ang{0.48}$; mean \omterm{def:b1:geological-maps:attitude}{dip direction}:
$2.94/\sqrt{12} = \ang{0.85}$.

\textbf{13.} No: \ang{2.9} of \omterm{def:b1:geological-maps:attitude}{dip direction} on a bed dipping \ang{12} is only
$2.9\sin\ang{12.1} = \ang{0.6}$ of turn of the plane, less than the
\ang{1.7} of the \omterm{def:b1:geological-maps:attitude}{dips}. Any slight undulation of a gentle bedding surface
scatters its \omterm{def:b1:geological-maps:attitude}{dip direction} widely.

\textbf{14.} $D = \ang{-2.5}$: $83.1 - 2.5 = \ang{080.6}$; her mean is
$080.6/12.1$, the same as the first geologist's.

\textbf{15.} The \omterm{def:b1:geological-maps:attitude}{dip} is read with the clinometer, against the vertical given
by gravity; the needle plays no part.

\textbf{16.} Yes: a rotation about the vertical moves every reading and the
mean by the same angle (\cref{exo:b1:field-lab-methods-1:11}); the answer is the same.

\textbf{17.} $10\times0.2 = \ang{2}$ of error on every \omterm{def:b1:geological-maps:attitude}{dip direction}, more
than twice the standard error of \ang{0.85}: a systematic error that the
averaging cannot detect.

\textbf{18.} $80.6 - 90 = \ang{350.6}$; the \omterm{def:b1:geological-maps:attitude}{attitude} is \omterm{def:b1:geological-maps:attitude}{strike} \ang{351},
\omterm{def:b1:geological-maps:attitude}{dip} \ang{12} to the east.

\textbf{19.} \omterm{def:b1:geological-maps:attitude}{Dip directions} $080.6$ and $080.6$: no difference at this
precision; \omterm{def:b1:geological-maps:attitude}{dips} $12.1$ and $12.2$: \ang{0.1}.

\textbf{20.} The normals $(\sin\ang{12.2}\sin\ang{80.6},
\sin\ang{12.2}\cos\ang{80.6}, \cos\ang{12.2})$ and $\vec S/\|\vec S\|$ make
an angle of \ang{0.13}.

\textbf{21.} $1.72/\sqrt{12} = \ang{0.50}$; the difference, \ang{0.13}, is a
quarter of it: not significant. The compass and the boreholes agree.

\textbf{22.} Four times as many: 48 readings.

\textbf{23.} For example: a wrong or outdated declination setting; a compass
deflected by iron nearby (a vehicle, tools, magnetic rocks); a clinometer
out of adjustment; readings taken only on the best exposed surfaces, which
may not represent the limb. Averaging reduces random errors, not systematic
ones.

\textbf{24.} The mean \omterm{def:b1:geological-maps:attitude}{attitude} of the east limb is \textbf{081/12}: \omterm{def:b1:geological-maps:attitude}{dip
direction} \ang{081}, \omterm{def:b1:geological-maps:attitude}{dip} \ang{12}.
\end{solution}
